\documentclass[11pt]{article}

\usepackage[a4paper,margin=1in]{geometry}
\usepackage{amsmath,amssymb,amsthm,mathtools}
\usepackage[T1]{fontenc}
\usepackage{lmodern}
\usepackage{microtype}
\usepackage{hyperref}

\newtheorem{theorem}{Theorem}
\newtheorem{lemma}[theorem]{Lemma}
\newtheorem{proposition}[theorem]{Proposition}
\newtheorem{corollary}[theorem]{Corollary}
\newtheorem{remark}[theorem]{Remark}
\newtheorem{definition}[theorem]{Definition}
\newtheorem{hypothesis}[theorem]{Hypothesis}
\newtheorem{conjecture}[theorem]{Conjecture}

\title{}
\author{}
\date{}

\begin{document}
\maketitle

For \(n\ge 2\) and \(0<a<1\), consider the real nonnormal tridiagonal Toeplitz matrix
\[
\mathbf{A}_n(a)\coloneqq \operatorname{tridiag}(a,0,1)
=
\begin{bmatrix}
0 & 1 \\
a & 0 & \ddots \\
& \ddots & \ddots & \ddots \\
&& \ddots & 0 & 1 \\
&&& a & 0
\end{bmatrix}\in \mathbb{R}^{n\times n}.
\]
For \(\varepsilon>0\), its \(\varepsilon\)-pseudospectrum is
\[
\sigma_\varepsilon\bigl(\mathbf{A}_n(a)\bigr)
\coloneqq
\left\{
z\in\mathbb C\colon
\bigl\|(z\mathbf I_n-\mathbf A_n(a))^{-1}\bigr\|_2>\varepsilon^{-1}
\right\},
\]
with the usual convention that the resolvent norm is \(+\infty\) on the spectrum \(\sigma\bigl(\mathbf{A}_n(a)\bigr)\).
Equivalently,
\[
\sigma_\varepsilon\bigl(\mathbf A_n(a)\bigr)
=
\left\{
z\in\mathbb C\colon
s_{\min}\bigl(z\mathbf I_n-\mathbf A_n(a)\bigr)<\varepsilon
\right\}.
\]
Then, there always exists a connectedness threshold
\[
N_{\text{c}}(a, \varepsilon)\coloneqq\min \left\{n \geq 2\colon \sigma_\varepsilon\bigl(\mathbf{A}_n(a)\bigr) \text { is connected}\right\}.
\]
And, with fixed \(a\) and \(\varepsilon\), for any \(n\ge N_{\text{c}}(a, \varepsilon)\), \(\sigma_\varepsilon\bigl(\mathbf{A}_n(a)\bigr)\) is connected.
Furthermore, if \(\sigma_\varepsilon\bigl(\mathbf{A}_n(a)\bigr)\) is connected, it is simply connected.


\end{document}